\documentclass[a4paper, 12pt, titlepage]{article}
\usepackage[utf8]{inputenc}
\usepackage[left=2cm,right=2cm,top=2cm,bottom=2cm]{geometry}
\usepackage[pdftex]{graphicx}
\usepackage[final]{pdfpages}
\usepackage[frenchb]{babel}%Pas important
\newcommand{\arc}[2]{$(#1~\rightarrow~#2)$}

\begin{document}
\newcommand{\myTitleOfTitlepage}{Séquenceur ADN}
\newcommand{\courseTitleOfTitlepage}{Algorithmique et bioinformatique}

\begin{titlepage}
\begin{center}
\includegraphics[scale=1.50]{figures/UMONS.jpg}\\[0.4cm]
\includegraphics[scale=0.30]{figures/FS_Logo.jpg}\\[3cm]
{\huge \courseTitleOfTitlepage}\\
\vspace{0.9mm}
\rule{9.5cm}{0.5mm}\\[0.5cm]
{\huge \bfseries \myTitleOfTitlepage}\\[0.2cm]
\rule{9.5cm}{0.5mm}\\
\textbf{Groupe 5}
\\[7cm]
      \begin{flushleft} \large
        Guillaume \textsc{Huysmans} 140223\\
        Jason \textsc{Bury} 130538\\
        Master 1 en Sciences informatiques\\
      \end{flushleft}
   \vfill
  {\large \today}
\end{center}
\end{titlepage}
\section{Introduction}
\subsection{Énoncé}
Le projet consiste à l'implémentation d'un séquenceur d'ADN,
c'est-à-dire reconstruire la séquence ADN cible à partir d'un ensemble de fragments.
Pour ce faire, une heuristique de type \emph{greedy} devra être utilisée pour
reconstruire un chemin hamiltonien (ce problème est NP-complet) dans un graphe
où les arcs sont des fragments et
les arcs sont pondérés avec les scores d'alignements semi-globaux.
Nous devons prendre en compte que certains fragments sont inclus à d'autre et
qu'il faudra parfois prendre le complémentaires inversé d'un fragment de
l'ensemble de départ.

\subsection{Consensus}
Les alignements deux à deux ne peuvent pas être utilisés tels quels :
d'une part, les gaps doivent parfois être reportés et
d'autre part, un vote doit être effectué lorsque des fragments se chevauchent.

\subsubsection{Report des gaps et ordonnancement}
Les fragments sont d'abord ordonnés par position croissante,
relative à celle du premier fragment.
Pour ce faire, le premier fragment est ajouté à la position 0 avant de
parcourir toutes les arcs \arc{s}{t} du chemin hamiltonien :
les gaps de l'alignement précédent sont ajoutés à celui de $t$ avec $s$ et
l'écart entre les alignements de $s$ et de $t$ (négatif si $s\subseteq t$) est
ajouté à la position courante. Cette position et ce nouvel alignement de $t$
sont ajoutés à une file de priorité.

Le calcul de l'alignement prend toujours un temps en $O(n^2)$, avec $n$ la
longueur du plus grand fragment. Le traitement dans \texttt{getAlignment}
est en $O(n)$ (parcours de toute la diagonale, dans le pire des cas).
L'ajout de $n$ éléments à une file de priorité se fait en $n\times O(1)=O(n)$.
Au total, on obtient du $O(n^2)$.

\subsubsection{Construction du contig}
Ensuite, on prend l'alignement au sommet de la file de priorité
(s'il est répété, on supprime ses répétitions) et on l'ajoute à
la liste des alignements en cours d'utilisation.
Ensuite, on itère tant qu'on n'a pas atteint le prochain élément de la file
(ou, si elle est vide, traité toute la liste de fragments) :
un vote est fait, les pointeurs sont avancés et les alignements terminés (ceux
dont le pointeur a dépassé la fin) sont retirés de la liste.
On recommence le processus jusqu'à ce que la file soit vide.

Le vote se fait par majorité et en cas d'égalité, $gap < A < C < G < T$.

Dans le pire des cas, on pourrait imaginer une suite de fragments préfixes les
uns des autres (l'entrée serait un genre de grand escalier descendant),
triés par ordre croissant de longueur. Le nombre de bases lues
serait en $O(n^2)$ avec $n$ la longueur du dernier fragment,
ce qui reste polynomial : au premier vote, on regarde $n$ nucléotides, au
second, $n-1$... et par la formule de Gauss, on trouve une complexité en
$n\times(n+1)/2=O(n^2)$.

En pratique, ils seront bien moins nombreux à être examinés en même temps.

\subsubsection*{Les inclusions de fragments}
Pour rappel, si nous avons deux fragments $f$ et $g$ et que lors de l'alignement, $f$ est inclus à $g$,
alors on a un arc \arc{f}{g} car le sens de l'arc signifie juste que l'alignement se termine à la fin de $f$.
Dans le calcul du consensus, à la première étape, celle qui consiste à calculer l'ordre de traitement des fragments,
nous recherchons pour chaque alignement, $\Delta$, la position du début de l'alignement à partir de la position du premier caractère de $f$.
Et cette position peut donc être négative dans le cas où $f$ est inclus à $g$ et que l'alignement commence avant le premier caractère de $f$.\\

En obtenant le $\Delta$ pour chaque alignement, nous pouvons donc obtenir une file de priorité contenant, dans l'ordre, les positions absolues de débuts d'alignements.
Par position absolue, nous entendons la position dans le contig sans suppression de gap.
Cette file de priorité nous premetra donc de savoir à partir de quelle position le nombre d'alignement à prendre en compte pour le consensus d'un symbole change.\\


Nous aurions pu nous simplifier la tâche et supprimer les fragments inclus à
d'autres mais leurs voix auraient été échangées contre un calcul plus simple
du consensus.
%TODO exemple simple ?


\section{Les complémentaires inversés dans l'algorithme Greedy}%explique comment on gère la suppressions des complémentaires inversés
Au cours de l'exécution de l'algorithme glouton, lorsqu'un arc est ajouté au chemin hamiltonien, les complémentaires inversés de ses noeuds ne doivent plus être pris en compte.
Pour cela, nous avons modifié l'algorithme glouton et ajouté une fonction dans la structure union-find:
en ajoutant un arc \arc{f}{g} au chemin hamiltonien, nous mettons les booléens \texttt{in} et \texttt{out} de leur complémentaire inversé à vrai.
Ainsi les arcs issus de ces noeuds ne seront jamais ajoutés dans le chemin.
De plus, nous avons ajouté une méthode dans l'objet \texttt{UnionFind} afin de décrémenter le compteur d'ensembles distincts dans la structure.
Cette méthode sera appelée deux fois : chaque fois qu'un arc \arc{f}{g} est ajouté au chemin.
En effet, si l'arc \arc{f}{g} est ajouté au chemin, nous savons que la structure \texttt{UnionFind} contient les ensembles \{$\bar{f}$\} et \{$\bar{g}$\} puisqu'ils n'ont jamais fusionné avec un autre ensemble.
Puisqu'ils ne seront plus jamais utilisés, il n'est pas nécessaire de
réellement les supprimer, il suffit de décrémenter de deux le nombre d'ensembles distincts.
Ce compteur d'ensembles distincts est utilisé pour arrêter l'algorithme
lorsqu'il n'y a plus qu'un ensemble dans l'\texttt{UnionFind}.

\section{Les inclusions de fragments}%TODO expliquer comment on gère les inclusions de fragment.

\section{Optimisations}
Dans cette section, nous expliquerons les techniques que nous avons utilisées dans le but d'améliorer les performances de notre programme, en temps de calcul ou en mémoire.

\subsection{Classe Complementary}
Puisque certains fragments peuvent être le complémentaire inversé d'un fragment
de la cible, nous devons créer et ajouter les complémentaires inversés de
tous les fragments dans le graphe avant l'application de l'algorithme Greedy.
Au lieu de créer ces fragments, nous associons à chaque séquence $f$ un objet
$\bar{f}$ qui hérite de \texttt{Sequence} et dont la méthode \texttt{get($i$)}
retournera le complémentaire du $i$\textsuperscript{ème} symbole
à partir de la fin de $f$.
Cela évite d'occuper deux fois la mémoire nécessaire
pour les fragments chargés à partir du fichier.

\subsection{Identifiant des fragments}
Pour améliorer les performances de l'algorithme Greedy, chaque fragment est
numéroté. Soient $n$ le nombre de fragments chargés depuis le fichier et $0 \le i < 2n$.
\begin{itemize}
 \item Si $i < n$, le fragment $i$ est le $(i+1)$\textsuperscript{ème} fragment chargé.
 \item Si $i \ge n$, le fragment $i$ est le complémentaire inversé du fragment $i-n$.
\end{itemize}

Les identifiants sont utilisés dans les arcs pour indicer la liste des fragments.

Cela améliore les performances. En effet, cela permet de remettre l'ensemble
des arcs du chemin hamiltonien dans l'ordre.
Lors de l'algorithme Greedy, à chaque fois qu'un arc \arc{f}{g} est ajouté
au chemin, on enregistre dans un tableau (indicé à partir de 0)
en position $f-1$ l'indice de l'arc dans le tableau où on vient de l'ajouter.
Une fois l'algorithme Greedy terminé, on trouve en temps linéaire
le premier fragment du chemin (c'est le seul dont le flag \texttt{in} est faux).
Ensuite, toujours en temps linéaire, on boucle en prenant chaque fois
l'arc \arc{f}{g} suivant que l'on trouve en position $f-1$ et
le prochain arc se trouve en position $g-1$.
Ainsi, remettre les arcs dans l'ordre se fait en $O(n)$
où $n$ est le nombre de fragments.
Si les noeuds n'étaient pas des identifiants, il aurait fallu effectuer
une recherche pour trouver l'arc suivant et
on aurait donc eu un algorithme en $O\left(n^2\right)$.

De plus, lorsqu'un arc \arc{f}{g} est ajouté au chemin hamiltonien,
trouver les complémentaires inversés de $f$ et $g$ pour mettre
leur \texttt{in} et \texttt{out} à vrai se fait en $O(1)$.
Ensuite, il n'y a plus besoin d'effectuer des recherches dans
la liste de fragments pour chaque appel à la méthode \texttt{find}
dans l'union-find si les noeuds sont représentés par un fragment.

Nous aurions pu nous passer d'identifiant en créant un objet \texttt{Node}
contenant un fragment, une référence vers le noeud du complémentaire inversé,
la position de l'arc suivant dans le chemin,
le parent dans l'union-find, \texttt{in} et \texttt{out} mais
cela rend l'implémentation plus compliquée et beaucoup moins élégante.

\subsection{Multithreading}
Avant de passer à l'algorithme Greedy, il faut créer tous les arcs et donc
calculer les scores d'alignement entre chaque paire de fragments.
Tous les calculs de score peuvent se faire en parallèle.
Deux méthodes de parallélisation sont envisageables:
\begin{enumerate}
\item Pour chaque thread, on définit une part égale de l'ensemble des arcs à calculer.
\item À chaque fois qu'un thread a calculé un arc, il pioche une autre tâche à
	effectuer dans une mémoire partagée avec les autres threads.
	C'est exactement le problème des producteurs/consommateurs. %FIXME formulation ?
\end{enumerate}
Nous avons choisi la première méthode car nous observons que les différences
de temps de calcul entre chaque thread sont identiques, à 1 seconde près,
%TODO refaire les tests et indiquer les résultats
ce qui signifie que leur charge de travail est équilibrée.
De plus, beaucoup d'appels systèmes auraient eu lieu si la deuxième méthodes
avait été utilisée.

\subsection{Même score pour plusieurs arcs}
Les arcs utilisés sont orientés. Le poids d'un arc \arc{f}{g} est
le score de l'alignement où soit on aligne le suffixe de $f$ au préfixe de $g$,
soit $f$ est inclus à $g$.
Après avoir calculé la matrice d'alignement, ce score est le maximum des scores sur la dernière ligne.
Il correspond aussi au score de \arc{\bar{g}}{\bar{f}} car l'alignement est exactement le même.
Le maximum des scores de la dernière colonne de la matrice donne celui de
\arc{g}{f} et donc aussi celui de \arc{\bar{f}}{\bar{g}}.
Quatre scores peuvent donc être calculés avec une seule matrice d'alignement.

\subsection{Typage des scores}
La valeur d'un score peut être représentée sur 2 octets (\texttt{short} en Java).
Effectivement, la taille maximale de fragment pour toutes les collections qui
nous ont été fournies est de 1366.
%TODO pas oublier de changer si changement des constantes de score
Une borne inférieure au score est -2732 (seulement des gaps)
et une borne supérieure est 1366 (alignement total).
Par conséquent, un \texttt{short} convient puisque sa valeur peut se trouver
entre -32768 et 32767.

\end{document}
